\documentclass[12pt]{article}
\usepackage{amsmath, amssymb}
\usepackage{geometry}
\geometry{margin=1in}

\begin{document}
	
	\title{Discussion on Niggli Reduction, $\varepsilon$-Modification, and Boundary Geometry}
	\author{User--Assistant Dialogue}
	\date{}
	\maketitle
	
	\section*{Conversation}
	
	\subsection*{User}
	This paper suggests modifying Niggli reduction by adding an ``epsilon'' to each of the tests for reduction steps.  
	Can you explain what adding epsilon to each of those tests changes the behavior at the edges of the Niggli cone?  
	Does it increase or decrease the Niggli cone at each boundary?
	
	\subsection*{Assistant}
	Adding an $\varepsilon$--tolerance to the inequalities in the Krivý--Gruber algorithm \emph{shrinks} the Niggli cone at every boundary.  
	It never enlarges it.  
	The effect is to push the algorithm away from boundary manifolds so that cells lying extremely close to a boundary are treated as if they were outside the Niggli region, forcing an additional reduction step.
	
	Mathematically, every Niggli inequality of the form
	
	
	\[
	x > y
	\]
	
	
	is replaced by something like
	
	
	\[
	x > y + \varepsilon,
	\]
	
	
	which removes a thin layer near the boundary.  
	Thus the admissible region becomes a strict subset:
	
	
	\[
	\text{Niggli}_\varepsilon \subset \text{Niggli}.
	\]
	
	
	
	This eliminates ambiguous cases, prevents cycling, and ensures termination.
	
	\subsection*{User}
	From their text for $x > y$, explain why $y < x - \varepsilon$ decreases the size of the cone and not increases it.
	
	\subsection*{Assistant}
	The original Niggli condition
	
	
	\[
	x > y
	\]
	
	
	defines the half-space $y < x$.  
	Replacing it with
	
	
	\[
	y < x - \varepsilon
	\]
	
	
	moves the boundary inward by $\varepsilon$.  
	Points with $x - y < \varepsilon$ that previously satisfied the inequality no longer do.  
	No new points are added.  
	Thus the region strictly shrinks.
	
	Geometrically, the boundary line $y = x$ is shifted to $y = x - \varepsilon$, reducing the allowed region.
	
	\subsection*{User}
	I want to know about each of those.
	
	\subsection*{Assistant}
	You asked about four topics:
	
	\begin{enumerate}
		\item The 45 Niggli boundary polytopes.
		\item Acute/obtuse ambiguity surfaces.
		\item The $G^6 \to S^6$ embedding and how $\varepsilon$ interacts with it.
		\item Perturbation stability near the cone edges.
	\end{enumerate}
	
	\paragraph{1. The 45 Gruber polytopes.}
	The Niggli cone in $G^6$ has 45 distinct boundary polytopes, each corresponding to one or more Niggli equalities becoming active:
	
	
	\[
	a^2 = b^2,\quad \xi = 0,\quad \xi + \eta + \zeta = 0,\quad \text{etc.}
	\]
	
	
	Adding $\varepsilon$ removes a small tubular neighborhood around each polytope, eliminating ambiguous reductions.
	
	\paragraph{2. Acute/obtuse ambiguity surfaces.}
	The signs of $\xi, \eta, \zeta$ encode whether angles are acute or obtuse.  
	The surfaces
	
	
	\[
	\xi = 0,\quad \eta = 0,\quad \zeta = 0
	\]
	
	
	separate acute from obtuse regimes.  
	Adding $\varepsilon$ excludes cells with nearly right angles, forcing the algorithm to choose a definite acute/obtuse regime.
	
	\paragraph{3. The $G^6 \to S^6$ embedding.}
	The Niggli cone is non-convex in $G^6$.  
	Your $S^6$ embedding straightens it.  
	The $\varepsilon$-shrinking in $G^6$ becomes a margin in $S^6$, giving a stable interior region for classification.
	
	\paragraph{4. Perturbation stability.}
	Without $\varepsilon$, the reduction map is multi-valued and discontinuous on boundaries.  
	With $\varepsilon$, the map becomes single-valued and stable on the interior:
	
	
	\[
	\|R(g + \delta) - R(g)\| \text{ is small for } \|\delta\| < \varepsilon.
	\]
	
	
	
	\subsection*{User}
	I have copied this chat, but can you render this whole chat as \LaTeX?
	
	\subsection*{Assistant}
	Yes.  
	This document is the \LaTeX{} rendering of the entire conversation.
	
	
	\subsection*{User}
	So it seems to me that what you are saying is that the epsilon removes from Niggli space all of the boundaries of the Niggli cone. Points that would fall there are assigned to some other place. In effect, the results of using non-zero epsilon are that points will be produced that do not meet the conditions for Niggli reduction. This seems to me to be an odd situation.
	
	\subsection*{Assistant}
	Your interpretation is correct. Introducing a non-zero $\varepsilon$ removes all boundary manifolds of the classical Niggli cone from the admissible region. Any point that would lie exactly on a Niggli boundary is treated as ``not sufficiently reduced'' and is pushed into the interior by an additional reduction step.
	
	This means that the $\varepsilon$-modified algorithm will never output a point that satisfies the classical boundary equalities such as
	
	
	\[
	A = B,\qquad \xi = 0,\qquad \xi + \eta + \zeta = 0.
	\]
	
	
	Instead, the algorithm produces a point strictly inside the cone, at least $\varepsilon$ away from every boundary.
	
	Although this may appear odd at first, it is intentional. The purpose of $\varepsilon$ is not to redefine the Niggli cone as a mathematical object, but to regularize the \emph{algorithm} so that it becomes stable, single-valued, and guaranteed to terminate. The classical Niggli boundaries are precisely where the reduction map is multi-valued, where acute/obtuse ambiguity occurs, and where cycling can happen. Removing a thin neighborhood around these boundaries eliminates all such pathological cases.
	
	Thus the $\varepsilon$-Niggli output is a point in the strict interior of the cone, not on its boundary. This ensures numerical robustness without altering the underlying geometric definition of the Niggli cone.
	
\end{document}
